\section{总结}\label{sec:summary}
在本工作中，我们利用赝 PDF 方法首次对 K 介子胶子部分子分布进行了格点研究。
我们以 clover 费米子为价作用，使用 MILC 合作组生成的 $\pi$ 介子质量 310~MeV、两个格距（0.15 与 0.12~fm）的 2+1+1 味 HISQ 规范组态。
我们采用动量涂抹与高统计测量（最多达 O(324,000) 次），使 K 介子的助推动量达到约 2~GeV。
我们利用两态拟合策略仔细研究了激发态对矩阵元的贡献，确保稳定地获得了基态矩阵元。
随后，我们用所得拟合基态矩阵元计算了约化赝 ITD（RpITD），并提取了胶子部分子分布。
我们发现，较细格距处的 K 介子胶子 PDF 在统计不确定度内与 DSE 结果~\cite{Cui:2020tdf} 一致，但小 $x$ 区除外——那里我们的 $z P_z$ 太小而无法约束。
在与 $\pi$ 介子 PDF 结果比较时，我们发现 K 介子 PDF 在大部分 $x > 0.2$ 区域中心值略小。
我们发现，K 介子胶子 PDF 在 0.15~fm 粗格距下显示出潜在的离散化误差；
未来补充 0.09~fm 格距的研究将使我们能更好地估计 0.12~fm 结果的系统不确定度。
根据此前关于 $\pi$ 介子胶子的计算，我们怀疑夸克—胶子混合小于我们的统计误差。
其他系统误差来源，如高扭度贡献、有限体积效应与非物理 $\pi$ 介子质量效应，未包含在这项开创性研究中。
未来研究应致力于改进这些系统误差，并为即将开展的实验工作提供对 K 介子胶子 PDF 更好的测定。
